\chapter{التطبيقات والتنقل}
\label{ch:builder-navigation}

\section{ماذا ستضبط في هذا الفصل؟}
\label{sec:builder-navigation-intro}

لا يكفي أن تنشر نوع السجلات حتى يستخدمه الفريق؛ فيجب أن يجده الأعضاء في المكان الذي يتوقعونه. ويشرح هذا الفصل صفحتين: «التطبيقات والتنقل» التي ترتّب التطبيقات وأقسامها في القائمة وتُظهرها أو تخفيها، و«التنقل حسب العمل» التي تحدد ما يراه كل عضو أولاً عند دخوله.

\begin{rcnote}[مثال: شركة الأفق]
بعد نشر «طلبات الزيارة» يضيف مسؤول النظام تطبيق «خدمات الفروع»، ويضع فيه قسماً يعرض الطلبات، فيجده عمر حداد في القائمة دون أن يبحث عنه.
\end{rcnote}

\section{صفحات هذا الفصل}
\label{sec:builder-navigation-pages}

الصفحتان مستقلتان في قسم «الأساسية» من أدوات البناء. ابدأ بالأولى لترتيب القائمة، ثم استخدم الثانية عندما تريد تجربة أوضح لكل دور.

\begin{table}[H]
\centering
\caption{إظهار العمل للفريق}
\label{tab:builder-pages-navigation}
\rcZebra
\begin{tabularx}{\linewidth}{R{4.2cm} Y}
\rcHeaderRow
\thh{الصفحة وقسمها} & \thh{ماذا تقدم هذه الصفحة؟}\\
التطبيقات والتنقل — الأساسية & ترتّب التطبيقات والأقسام التي يراها الأعضاء، وتُظهرها أو تخفيها.\\
التنقل حسب العمل — الأساسية & تحدد ما يراه كل عضو أولاً: صفحة «مكتبي» وقائمة المستحق عليه والسياق المشترك.\\
\end{tabularx}
\end{table}
